% D5 finite-cover normal form 044: theorem and lemma snippet

\begin{theorem}[Checked finite-cover normal form for the best-seed D5 active branch]
Fix the best seed pair
\[
L=[2,2,1],\qquad R=[1,4,4],
\]
and let $A_m$ be the extracted active branch of the unresolved channel
\[
R1 \to H_{L1}
\]
for a checked modulus $m\in\{5,7,9,11\}$. For $x\in A_m$, define
\[
B(x)=(s(x),u(x),v(x),\lambda(x),f(x)),
\qquad
c(x)=\mathbf 1_{\{q(x)=m-1\}},
\]
where $\lambda=\mathrm{layer}$ and $f\in\{\mathrm{regular},\mathrm{exceptional}\}$ is the carried family bit. Let
\[
\tau(x)=\min\{t\ge 0: c(F_m^t x)=1\},
\qquad
U^+(x)=u(F_m^{\tau(x)}x),
\]
and set
\[
d(x)=\mathbf 1_{\{U^+(x)\ge m-3\}}.
\]
Then the checked active branch admits the normal form
\[
B \leftarrow B+c \leftarrow B+c+d,
\]
with the following properties:
\begin{enumerate}
\item the exceptional fire predicate is a function of $B$;
\item the regular fire predicate is not a function of $B$, but it is a function of $B+c$;
\item every $(B+c)$-fiber with $c=1$ is a singleton;
\item every non-singleton $(B+c)$-fiber lies on the regular noncarry region;
\item the checked residual $2$-sheet cover over $B+c$ is realized by the binary anticipation bit $d$;
\item the active transition law is deterministic on $B+c+d$.
\end{enumerate}
\end{theorem}

\begin{lemma}[Future-carry low/high dichotomy on the ambiguous support]
For each checked modulus $m\in\{5,7,9,11\}$, if $x$ lies in a non-singleton $(B+c)$-fiber, then
\[
U^+(x)\in\{0,1,m-3,m-2\},
\]
with the expected degeneracy at $m=5$. In particular, the threshold bit
\[
\mathbf 1_{\{U^+(x)\ge m-3\}}
\]
produces a binary separation of every checked ambiguous $(B+c)$-fiber.
\end{lemma}

\begin{corollary}[Trigger-level vs structure-level lift]
On the checked active branch, the exact trigger-level missing lift is
\[
c=\mathbf 1_{\{q=m-1\}},
\]
while the residual structure-level lift is the binary anticipation sheet
\[
d=\mathbf 1_{\{U^+\ge m-3\}}.
\]
Thus the local realization frontier remains the carry sheet, not the full residual cover.
\end{corollary}
